\chapter{Smooth Necks, Caps, and Their Global Topology}
\label{ch:neck-cap-topology}

\section{Purpose and interfaces}

This chapter supplies the topological part of the canonical-neighborhood
argument.  Its input is a three-dimensional ancient $\kappa$-solution from
Chapter~\ref{ch:ancient-models}, or a fixed Riemannian three-manifold carrying
a sufficiently fine neck--cap cover.  Its output is a genuine smooth model:
an open cylinder, a cap attached to one end, a tube with one or two caps, a
circle bundle of spheres, or a closed component modelled on $S^3$,
$\mathbb{RP}^3$, or $\mathbb{RP}^3\#\mathbb{RP}^3$.  The
distinction between a local cover and a cover of the whole connected component
is essential.  The surgery chapter uses the smooth ball and sphere data; the
singular-limits chapter uses only the local neck and cap alternatives.

The models retain their actual embeddings, attaching collars, and central
spheres.  These data are needed when a later surgery uses a particular
sphere inside the original manifold.  The sphere filling and isotopy inputs
are described in Section~\ref{sec:ball-gluing}.

\section{Necks, caps, and covers}
\label{sec:neck-cap-definitions}

Let $M$ be a $T_3$ smooth manifold modelled on $\mathbb R^3$, with Borel
measurable structure, and let $g$ be a Riemannian metric.  An
\emph{$\varepsilon$-neck} consists of $0<\varepsilon<1/2$, a centre $p$,
and a scale $r>0$ with
\[
  r=R(p)^{-1/2}>0,
\]
together with a smooth coordinate diffeomorphism
\[
 \Phi:S^2\times(-\varepsilon^{-1},\varepsilon^{-1})\longrightarrow U\subset M.
\]
The central sphere is $\Phi(S^2\times\{0\})$.  The coordinate metric is
required to be $\varepsilon$-close to the round cylinder after rescaling by
$r^{-2}$.  Precisely, put $g_{\rm cyl}=2g_{S^2}+ds^2$ and
$m=\lfloor\varepsilon^{-1}\rfloor$.  The rescaled tensor is smooth, and
there is a single $b<\varepsilon^2$ such that, at every point of the domain,
\[
 \sum_{j=0}^{m}\bigl|\nabla_{g_{\rm cyl}}^j
       (r^{-2}\Phi^*g-g_{\rm cyl})\bigr|_{g_{\rm cyl}}^2\le b.
\]
For $a<b$ write
\[
 U(a,b)=\{x\in U:a<\operatorname{pr}_2\Phi^{-1}(x)<b\}.
\]
The source uses these regions, rather than an informal ``left'' and ``right''
side, in every overlap statement.  Reversing the axial coordinate is allowed;
two necks are equivalent up to reversal when their scales, centres, carriers,
central spheres, and coordinates agree after $s\mapsto\sigma s$ with
$\sigma\in\{1,-1\}$.

An \emph{$\varepsilon$-neck-only cover} is a connected set $X\subset M$ and a
set of necks such that every $x\in X$ is the centre of one of them, all have
the same $\varepsilon$, and $\varepsilon$ is below a threshold at most
$1/200$.  A \emph{connected neck--cap cover} additionally contains caps whose
cores cover the points not selected as neck centres.  A cover is \emph{whole}
when $X=M$; no theorem about a local cover may silently replace $X$ by $M$.
Each cap has an open carrier, a nonempty core equal to the interior of a
compact closed core, a smooth end neck, a boundary neck and boundary sphere,
and a model equivalence with either Euclidean space or a punctured real
projective three-space.  The connection used by the cap, its end neck, and its
boundary neck is the same connection.  The cap estimates include positive
scalar curvature, an intrinsic-diameter bound, a scalar-ratio bound, a volume
bound, and gradient and scalar-evolution bounds.  For every core point $y$ a
positive radius $\rho(y)$ gives a compactly contained metric ball with
\[
 \sup_{B(y,\rho(y))}R=\rho(y)^{-2}
\]
and a lower volume bound proportional to $\rho(y)^3$.

Here are the quantitative requirements more explicitly.  Write $V$ for the
cap carrier, $Q$ for its closed core, $c>0$ for its cap constant, and
$R_V=\sup_V R$.  Then $Q=V\setminus U_{\rm end}$,
$\operatorname{int}Q$ is the core, and
\[
 \operatorname{diam}_{V}V<cR_V^{-1/2},\qquad
 \operatorname{Vol}(V)<cR_V^{-3/2}.
\]
The diameter uses the infimum of lengths of smooth paths constrained to $V$.
There exist constants $b_R,b_\nabla,b_\Delta<c$ such that, for all $x,y\in V$,
\[
 R(y)\le b_RR(x),\quad |\nabla R|(x)\le b_\nabla R(x)^{3/2},\quad
 |\Delta R+2|\operatorname{Ric}|^2|(x)\le b_\Delta R(x)^2.
\]
There is also one $b_v>c^{-1}$ for which every core ball above has volume at
least $b_v\rho(y)^3$.  These auxiliary constants are chosen for the cap,
before the points.  The boundary is both $\partial Q$ and the central sphere
of the boundary neck.  Near every boundary point it is $f=0$, where $f$ is
smooth, $df\ne0$, and $Q$ is locally $f\le0$; it lies in the closure of the
negative terminal quarter of the end neck.  Thus the data specify a smooth
domain and an attaching collar, as well as estimates.  The same scale and
Levi--Civita connection are retained when the cap is glued to a tube.

\subsection{Scalar control on the entire neck}
\label{sec:neck-cap-full-scalar}

The neck estimates hold up to every point of the open coordinate domain;
the distance of a point from its ends need not have a fixed positive lower
bound.  The quantitative estimate needed in the singular-limits chapter is
\begin{equation}
 |r^2R(x)-1|\le 10^{12}\varepsilon
 \qquad(x\in U,\quad 0<\varepsilon\le10^{-6}).
 \label{eq:neck-cap-full-scalar}
\end{equation}
Both numerical constants precede the choices of manifold, metric, neck and
point.  The scalar may be computed using any Levi--Civita description of
the same metric, since that connection is unique.

Here is the two-jet calculation giving this uniformity.  At any cylinder
point $(q,s)$ use a rotated sphere chart and translate the axial coordinate
so that the point becomes zero.  A neighborhood of zero fits inside the
neck even when $s$ is near an end; its radius can depend on $s$.  Only the
metric germ at zero enters the calculation.  In these coordinates the
normalized metric and cylindrical metric have two-jets $J,J_0$ with
zeroth, first and second derivative errors bounded by
$\delta=19044\varepsilon\le1/2$.  Indeed, each scalar coefficient and
its first two derivatives have error at most $2116\varepsilon$, and
the nine coefficients give this tensor bound.  The model first derivative
vanishes, its second derivative has norm at most $1000$, and its metric
and inverse have norms at most $2$ and $1$.  Quantitative inversion gives
\[
 \|J^{(0)}\|\le3,\qquad
 \|(J^{(0)})^{-1}\|\le2,\qquad
 \|(J^{(0)})^{-1}-(J_0^{(0)})^{-1}\|\le2\delta.
\]
The coordinate curvature formula consists of terms linear in the second
metric derivatives and terms quadratic in the first derivatives, contracted
with inverse metrics.  Subtract these expressions for $J$ and $J_0$ one
factor at a time.  The above bounds, followed by the two scalar contractions,
give
\[
 |\operatorname{Scal}(J)-\operatorname{Scal}(J_0)|
 \le(20088+972\cdot1000)\delta<10^{12}\varepsilon.
\]
The model scalar is one, and naturality and constant rescaling identify the
other scalar with $r^2R(x)$.  This proves
\eqref{eq:neck-cap-full-scalar} on the full carrier.  It is a numerical
version of the neck curvature comparison in
\cite[Lemma A.2, pp.~497--498; Proposition A.11, pp.~503--504]{MT2007}.
Since $r^2R(p)=1$, further reducing $\varepsilon$ so that
$10^{12}\varepsilon<1/2$ gives
$R(p)/2<R(x)<2R(p)$ everywhere in the neck.  The linear estimate, rather
than only this factor-two consequence, is what controls a change to the
terminal scalar normalization in Chapter~\ref{ch:singular-limits-v4}.

\section{Overlap geometry and balanced tubes}
\label{sec:balanced-tubes}

Two consecutive necks in a balanced chain overlap in the positive half of one
and the negative half of the other.  More precisely, if $N_i,N_{i+1}$ have
parameter $\varepsilon$, then
\[
 N_i(\varepsilon^{-1}/2,\varepsilon^{-1})\subset N_{i+1},\qquad
 N_{i+1}(-\varepsilon^{-1},-\varepsilon^{-1}/2)\subset N_i,
\]
while their intersection lies in the corresponding three-quarter slabs.
The centres are distinct, and the distance between adjacent centres lies in
the interval $[.99r_i\varepsilon^{-1},1.01r_i\varepsilon^{-1}]$ (with the
extended metric interpretation).  For every pair of active indices $i<j$,
there exists $s\in(-\varepsilon^{-1},0)$ such that
$U_j\cap U_i(-\varepsilon^{-1},s)=\varnothing$.  A chain may be finite,
one-sided infinite, or bi-infinite; its active index set is respectively an
integer interval, a half-line, or all of $\mathbb Z$.

The first global construction is a maximal-chain argument.  Start with a
selected neck and extend at a positive frontier whenever the next covered
centre is available; prepend at a negative frontier in the analogous way.
Quarter-capture says that a missed point would produce a new source neck at a
finite endpoint.  The no-return estimate supplies a negative interval on
which the new carrier is disjoint from the old negative end.  Either extension
would contradict maximality.  Hence a separating neck-only cover is contained
in the union of one balanced chain.  The chain's overlap charts paste to an
open-cylinder model, and the central spheres are smoothly isotopic inside that
model.  This is Proposition~A.19, not a mere covering-space slogan.

The separation label is component-relative:
\[
 \operatorname{Sep}(N)\iff
 (\pi_0(N_\mathrm{centre})\setminus S_N)\ne\varnothing
 \text{ and is disconnected},
\]
where $S_N$ is the central sphere.  The complementary connectedness label is
\(\operatorname{NonSep}(N)\) when that complement is connected.  A positive
axial slice proves that the complement is nonempty, so the labels are
exclusive and exhaustive.  Nearby centred necks have ambient-isotopic central
spheres; therefore the label is locally constant on a connected centre cover.
This is the corrected form of the source's separation convention and is why
unrelated ambient components cannot prove separation.

\section{The nonseparating branch}
\label{sec:circle-bundle}

Suppose a whole neck-only cover has only nonseparating necks.  We first
construct a finite returning chain; compactness of $M$ is a consequence.
Fix a source neck $N$ and put $L=\varepsilon^{-1}$.  Its central sphere
has connected complement.  That complement is an open manifold and hence
locally path connected, so a path in it joins
$N(q,-L/2)$ to $N(q,L/2)$.  Let $K$ be the compact image of this path
together with the original centre.  At an outgoing frontier the
extension-or-return construction selects a source neck centred on the path,
outside every earlier carrier.  Either it extends the balanced chain or
its carrier meets every negative tail $N(S^2\times(-L,c))$, $-L<c<0$.
The latter is the deep return used below.

There is a uniform finite bound on the number of extensions along this
fixed path.  A finite source-neck subcover of $K$, and scale comparison
at an intersecting pair of necks, bound all candidate scales below by
some $r_*>0$.  The axial-depth estimate then separates any earlier and
later centres by at least
$d_*=r_*\sqrt{1-\varepsilon}L>0$, since the later centre is outside the
earlier carrier.  A finite cover of $K$ by metric balls of radius
$d_*/3$ contains at most one of these centres in each ball.  An
arbitrarily long chain is therefore impossible, and a finite deep return
occurs.  This uses only compactness of the chosen path.

The terminal and initial slices of this chain, together with
the return neck, cut the union into three compact pieces.  The closing
construction gives smooth product charts on the three pieces.  Their common
boundaries are exactly the lower boundary spheres of the next piece, so the
charts paste to a smooth projection
\[
                 p:M\longrightarrow S^1.
\]
Each fibre is a smoothly embedded two-sphere, and every neck's central sphere
is smoothly isotopic in $M$ to the fibre over the corresponding circle point.
The local trivialisations preserve the projection.  Thus $M$ is represented
by a smooth $S^2$-bundle over $S^1$, with the possible monodromy retained in
the bundle model.  The certificate includes a homeomorphism to its model,
smooth forward and inverse maps, compactness, and the equality of every neck
accuracy parameter with the cover's $\varepsilon$.

Here is the step that ensures that the return covers the whole manifold.
Let $U$ be the union of the finite chain, let $R$ be the return neck, and
let $Q$ be the auxiliary closing neck.  Cutting off the two outgoing parts of
$U$ leaves a compact region $K$.  The two closing strips $D_R,D_Q$ have
compact closures, and the overlap calculation gives
\[
 W=(U\cup R)\cup Q=K\cup\overline{D_R}\cup\overline{D_Q}.
\]
Consequently $W$ is compact.  It is open because it is a union of open neck
carriers, closed because the ambient manifold is Hausdorff, and nonempty
because it contains the initial neck centre.  The whole-cover hypothesis
makes $M$ connected, so $W=M$.  This is the precise use of connectedness
that prevents a returning chain from describing only part of the manifold.

The three product charts $e_i:S^2\times[0,1]\longrightarrow M$,
$i\in\mathbb Z/3\mathbb Z$, have images respectively
$K,\overline{D_R},\overline{D_Q}$ and extend smoothly over collars of their
boundaries.  The intersection of consecutive images is exactly
$e_{i+1}(S^2\times\{0\})$.  After a fibrewise change of the interval
coordinates, the circle coordinate is the interval coordinate on the
$i$th of three successive arcs of $S^1$.
The collar matching makes these coordinates agree at each endpoint and on
its neighborhood.  They therefore define a smooth map, and the same
extended charts are its local trivialisations.  Smooth isotopies first move
each chain sphere to the terminal slice, and then to a fibre in these
charts.  The transition around the last arc can act nontrivially on $S^2$;
no triviality of monodromy is asserted.

This proof uses a finite return, not an assumed periodic parametrisation.  The
three-piece decomposition is what supplies the global projection; the
connectedness of the cover is what lets the local separation labels propagate
around the entire circle.

The overlap and smooth matching deserve a more explicit description.
Orient the return neck as $\widehat R$, and choose a source neck, possibly
reversed, $Q$ centred at height $-17L/20$ in $N$.  The overlap estimates
give smooth transition maps from the two bands of $Q$
\[
 S^2\times(-L/40,L/40),\qquad
 S^2\times(-13L/50,-6L/25)
\]
into $N$ and $\widehat R$, respectively.  Their axial derivatives are
positive.  The terminal slice of the chain is the graph $s=f(q)$ in
$\widehat R$, the zero slice of $Q$ is $s=h_N(q)$ in $N$, and the
$-L/4$ slice of $Q$ is $s=h_R(q)$ in $\widehat R$, with
\[
 -3L/10<f<-L/5,\quad -9L/10<h_N<-4L/5,
 \quad -L/5<h_R<19L/20.
\]
In particular $f<h_R$.  Cutting between these spheres leaves the
compact middle of the chain, the graph strip $f\le s\le h_R$ in
$\widehat R$, and the strip $-L/4\le s\le0$ in $Q$.  The end
exclusions of the balanced chain identify their pairwise intersections
as precisely the three indicated boundary spheres.  No assertion about
the entire intersection of two uncut neck carriers is needed.

For a graph strip $a(q)\le s\le b(q)$, its product chart is given by
$s=(1-u)a(q)+ub(q)$.  Positivity of $b-a$ and compactness of the
sphere extend this chart to a common interval $-\eta<u<1+\eta$.
The finite chain's cylinder chart gives the same construction for its
middle piece.  These are the charts $e_i$ above.  Matching their boundary
sets alone would not match their derivatives.  In the outgoing collar of
$e_i$ the required scalar coordinate is
\[
 g_i(q,u)=1+\operatorname{pr}_2(e_{i+1}^{-1}(e_i(q,u))).
\]
It equals one at $u=1$ and has positive $u$ derivative.  Extend the two
positive endpoint germs $u$ and $g_i$ to a strictly increasing smooth
coordinate $d_i(q,u)$ on the whole interval, fixing zero and one.
Here is the positive splice used in this construction.  Shrink the lower
collar until its germ $g_0$ is less than $1/4$.  Blend it with
$\ell(u)=3/8+u/4$ using an increasing cutoff $\beta$ which changes
from zero to one inside that collar.  The derivative of the blend is
\[
 (1-\beta)\partial_u g_0+\beta/4+\beta'(\ell-g_0)>0.
\]
The value gap makes the last term nonnegative.  Apply the same construction
to $1-g_i(q,1-u)$ at the upper end and reflect it back.  Both splices
equal $\ell$ in the middle, since $1-\ell(1-u)=\ell(u)$, so they
paste smoothly.  Extending the positive coordinate to the line gives
the required fibre diffeomorphisms; their smooth inverses follow from
the parameter-dependent inverse function theorem.
On a common raw collar the corrected coordinates satisfy
$d_i(q,s)=s$ and
$d_{i-1}(e_{i-1}^{-1}e_i(q,s))=1+s$.
Thus the formula
\[
 p(e_i(q,u))=\exp\bigl(2\pi\mathrm{i}(i+d_i(q,u))/3\bigr)
\]
agrees smoothly across each join, including the last join modulo three.
The corrected product charts over small circle arcs have source exactly
$p^{-1}(V)$, which proves local triviality over the entire fibre, not just
near one point of it.  This supplies the constructive form of
\cite[Lemma A.20, p.~508]{MT2007}; the separating and nonseparating cases
in the first two sentences of its printed proof are interchanged.

\section{Caps, encounters, and sphere gluing}
\label{sec:ball-gluing}

For a connected cover containing a cap core, follow the outward end of that cap
through successive necks.  At each step either the next cap core is disjoint,
or it meets the current frontier and can be enlarged.  A uniform cap bound
turns each frontier encounter into a definite increase in the distance from a
fixed point to the complement.  Since the cap diameter supplies a common upper
bound, an infinite nested sequence is impossible.  Well-foundedness therefore
produces a maximal cap with no admissible frontier enlargement.

The quantitative termination argument is useful because neither compactness
of the ambient manifold nor a bound on the number of necks is assumed.
There is a universal $0<\varepsilon_0\le1/200$ such that the following holds
for every $B$, every $0<\varepsilon\le\varepsilon_0$, and every sequence of
caps $V_n$ with this accuracy and cap constants at most $B$.  It is impossible
to have, for every $n$,
\[
 V_n\subset V_{n+1},\qquad
 \partial V_n\cap Q_{n+1}\ne\varnothing.
\]
To prove this, choose $p$ in the first core.  Nesting puts $p$ in every $V_n$,
and the scalar curvature at $p$ is independent of the presentation of the
Levi--Civita connection.  Put
\[
 \delta=.38(BR(p))^{-1/2}\varepsilon^{-1}>0,\quad
 L=BR(p)^{-1/2},\quad
 a_n=\inf_{z\notin V_n}d_g(p,z).
\]
The nested-frontier estimate gives
$a_n+.38r_{n+1}\varepsilon^{-1}\le a_{n+1}$, where $r_{n+1}$ is the
boundary-neck scale of the next cap.  The scalar-ratio bound gives
$r_{n+1}\ge(BR(p))^{-1/2}$, so $a_n+\delta\le a_{n+1}$.
To obtain the opposite bound, take a point of
$\partial V_n\cap Q_{n+1}$.  Since $V_n$ is open, that point is in its
complement, and the cap diameter estimate, extended to the closure, gives
$a_n\le L$.  Induction gives $n\delta\le a_n\le L$.  Choosing an integer
$n>L/\delta$ is a contradiction.  The argument is valid for extended metric
distances because the upper bound $L$ is finite.  The positive number $.38$
is the margin retained in the preceding overlap estimate; no claim about an
optimal margin is needed.  Applying this absence of infinite growth to the
strict cap-enlargement relation yields the maximal cap used above.

For completeness, the distance-increment estimate has two geometric cases.
If the earlier carrier misses the later boundary sphere, the disjoint-boundary
lemma puts the earlier carrier in the later core.  Crossing from that core to
the complement of the later cap costs at least $.9r\varepsilon^{-1}$, so the
weaker $.38r\varepsilon^{-1}$ estimate follows.  Otherwise the later boundary
meets the earlier carrier.  Frontier contact excludes containment of the
entire later boundary in the earlier carrier; hence this is a mixed-boundary
configuration.  The mixed-boundary neck estimate separates every point of
the earlier frontier from every point of the later complement by
$.38r\varepsilon^{-1}$.  Every path from $p$ to the latter complement first
crosses the earlier frontier.  Splitting its length at that crossing and
taking infima proves the required additive gain.  Both cases use a threshold
chosen before the cap constants; their minimum gives the uniform threshold
in the termination statement.

To see the numerical margin in the mixed case, write $C$ for the old cap,
$D$ for the new cap, $r_C$ for the old end-neck scale, and $r_D$ for the
new boundary-neck scale, and put $\ell=\varepsilon^{-1}$.  The connected
boundary sphere of $D$ meets
$C$ but is not contained in it, so it crosses $\partial C$ at a point
$x$.  The whole frontier $\partial C$ lies in the closure of the old
positive quarter.  For every $z\in\partial C$ and $y\notin D$ the
quarter-diameter and boundary-clearance estimates give
\[
 d_g(x,z)\le .51r_C\ell,\qquad d_g(x,y)\ge .9r_D\ell.
\]
Scale comparison holds at this common closure point by continuity of
scalar curvature and gives $r_C\le1.01r_D$.  The triangle inequality
therefore yields
\[
 d_g(z,y)\ge(.9-.51\cdot1.01)r_D\ell
             =.3849r_D\ell\ge.38r_D\ell.
\]
For extended distances the same argument cancels the finite term
$.51r_C\ell$ from an additive inequality.  It requires neither a minimizing
geodesic nor a completeness assumption.  This is the clearance used in
the growth argument, refining \cite[Claim A.22, pp.~509--510]{MT2007}.

The finite stopping tuple is more precise than this summary.  It contains a
finite balanced chain beginning with the original cap's end neck; active necks
are separating, their positive centres lie outside the original carrier, and
adjacent quarter slabs are mutually contained.  The chain forms a capped tube
whose carrier is the union of the cap and all active neck carriers.  If a
second cap core remains outside that carrier, the second-overlap analysis gives
one of two outcomes.  Either the two cap carriers form a compact connected
closed component, or their boundaries admit a second cylinder and the two
capped tube is assembled.  The two-cap and doubly-capped cases use different
overlap hypotheses; they are not interchangeable.

The compact closed-component step has two inputs.  A smooth Schoenflies
service says that a smooth embedded two-sphere with its collar extends to a
ball chart in Euclidean three-space.  Its proof fills the bounded side and
then extends the resulting ball embedding to an ambient diffeomorphism.  A
smooth sphere-isotopy service, obtained from Smale's sphere theorem and
Munkres' compactly supported plane isotopy, makes the boundary parametrisations
agree.  For two Euclidean caps their union is a three-sphere; for a Euclidean
cap and a punctured-projective cap it is a real projective three-space.  The
two punctured-projective case gives
$\mathbb{RP}^3\#\mathbb{RP}^3$ by the separate projective-pair producer.  The
result specifies the model of the closed component; it makes no assertion
that every component of the ambient manifold is a sphere.

\subsection{Filling the bounded side and extending its ball chart}
\label{sec:neck-cap-ambient-filling}

We distinguish the construction of a ball from the extension of an already
parametrized ball.  For a smooth embedded sphere in $\mathbb R^3$, smooth
Schoenflies states that its bounded complementary side is a smooth ball
\cite[Theorem 1.1, pp.~1--5]{Hatcher3M2014}.  The proof used here follows
the smooth Morse-cut argument, including its reflected saddle cases.
Perturb the height by a small supported ambient isotopy to make its
critical points nondegenerate with distinct heights.  Choose regular
levels between them.  At each such level, the finitely many section
circles bound planar disks.  Cut first along an innermost disk, round
the two resulting caps, and repeat.  Since every circle on the sphere
separates, each cut produces two spheres.  Keeping the actual cap disks
and the unchanged embedding near each original critical point gives a
finite tree of surgeries whose terminal cores have at most one original
critical point.

On a regular band, the vector field
$\nabla h/|\nabla h|^2$ identifies the height levels.  It gives the
cylindrical terminal model when there is no critical point; the Morse
charts give the rounded minimum and maximum models.  For a saddle the
chart has height $c-u^2+v^2$.  Its terminal core has three annular ends.
There are either two lower ends and one upper end, or the reverse.
Indeed the two possible height orientations cannot both have one lower
end: regular-band transport would then account for only two of the three
ends.  In the first orientation the two lower planar circles are either
disjoint with disjoint interiors or nested.  These give the two saddle
models; height reflection gives the other two.  This is the saddle part
of \cite[Lemma 1.2, pp.~3--4]{Hatcher3M2014}.

The matching is relative to a protected neighborhood of the Morse chart.
Straighten the regular annuli in the height direction, and match their
planar circles and connecting strips to the disjoint or nested model.
The planar isotopies extend with compact support, fixed on the protected
chart.  Their smooth dependence on height gives an ambient map
$(x,z)\mapsto(\Phi_{\chi(z),z}(x),z)$, where $\chi=1$ on the retained
band.  Its inverse uses the inverse planar family at the same height.
The three remaining caps are matched simultaneously, with their original
labels, by the relative disk replacement construction: choose common
closing disks, use the disjoint-or-nested order to make innermost
replacements first, and keep the already matched band fixed.  Arbitrary
independent cap maps would not ensure that their supports respect the
other caps.  For the reverse height orientation, conjugate the planar
family and the single cap-replacement map by
$(x,z)\mapsto(x,-z)$ and reverse the interval endpoints.  This preserves
the whole labelled caps, their disjointness, and the protected region.
The terminal sphere is thereby the ambient image of the boundary of its
filled standard model.

Reverse the finite surgery tree.  When two filled spheres are joined
back along their marked disk, their balls are either disjoint apart
from that disk, or one is nested inside the other.  The restored region
is respectively a ball with a ball attached, or a ball with the nested
ball removed along the marked boundary disk.  Both are diffeomorphic
to a ball: normalize the marked boundary disk and push its boundary
across the attached ball, using isotopy extension, as in
\cite[Lemma 1.3, p.~5]{Hatcher3M2014}.  Concretely, choose the child
ball whose filling avoids the other child's boundary, replace it by a
thin lens while fixing the surgery annulus and the other child, and
stretch the other child's cap across that annulus.  The two resulting
boundary images agree exactly.  Undoing the lens replacement and the
initial straightening transports the remaining filled ball to the
original sphere.  The rounded collars retained
at each cut make these smooth attachments.  Induction and the inverse
of the initial perturbation fill the original sphere.  The bounded-side
identification then ensures that the closed-ball image is the prescribed
closed domain, not merely some region with diffeomorphic boundary.

For a given full collar $c:S^2\times(-\delta,\delta)\to\mathbb R^3$,
separation supplies two complementary open regions.  Exactly one is
bounded; reverse the collar coordinate if necessary so its closure $K$
is locally $s\le0$.  Thus $K$ is compact,
$\overline{\operatorname{int}K}=K$, and its entire frontier is the
central sphere.  Collar charts give smooth half-space charts on $K$.
Its interior is connected: its intersection with the collar is the
connected inward half, while any additional component would have a
frontier point on that same collar.  The filling just proved therefore
gives a ball chart with closed-ball image exactly $K$.  A sphere
reparametrisation, extended across the ball using the isotopy below,
and a supported collar correction make this chart agree with
$c(q,s)$ at $\exp(s)q$ on a small collar of the unit sphere.

Finally let $b$ be any smooth ball chart on a neighborhood of
$\overline B_1$, injective there with invertible derivative.  A smooth
cutoff extends its germ to a globally smooth map $f$ agreeing near
$\overline B_1$; global injectivity of this auxiliary map is unnecessary.
For $0<a\le1$, the embeddings
\[
 x\longmapsto f((1-t+ta)x),\qquad 0\le t\le1,\quad |x|\le1,
\]
have invertible derivatives.  Extend their velocity by a compactly
supported vector field in ambient space and integrate it.  The resulting
diffeomorphism $P$ satisfies $P(f(x))=f(ax)$ on $\overline B_1$.
Put $A=Df(0)$.  For sufficiently small $a$, the local interpolation
between $f(y)-f(0)$ and $Ay$ has invertible derivative on $|y|\le a$,
uniformly in time.  The same velocity extension gives a supported
diffeomorphism $\Psi$ with
$\Psi(f(y)-f(0))=Ay$ there.  Consequently the ambient diffeomorphism
\[
 Q=a^{-1}A^{-1}\circ\Psi\circ (x\mapsto x-f(0))\circ P
\]
satisfies $Q(f(x))=x$ on $\overline B_1$.  Its inverse extends $b$
on that closed ball, including its actual sphere parametrisation.
This argument starts with a filled ball and introduces no sphere-filling
assumption into the preceding Morse construction.

\subsection{Prescribed collars and sphere parametrisations}
\label{sec:neck-cap-prescribed-collars}

The collar extension retains more information than the diffeomorphism type
of the filled ball.
\begin{lemma}[A ball chart with the prescribed outer collar]
\label{lem:neck-cap-prescribed-collar}
Suppose
$\psi:S^2\times(-1,1)\longrightarrow\mathbb R^3$ is a smooth collar
embedding.  There are a sign $\sigma\in\{1,-1\}$, a bounded connected
open set $\Omega$ with boundary $\psi(S^2\times\{0\})$, a sphere
diffeomorphism $h$, and a diffeomorphism $J$ from $B(0,2)$ onto
\[
 \Omega\cup\psi(S^2\times\sigma[0,1))
\]
such that $\psi(S^2\times\sigma(-1,0))\subset\Omega$, the complement
of $\overline\Omega$ is connected, and
\[
 J((1+s)h(q))=\psi(q,\sigma s)\qquad(0\le s<1).
\]
In particular, for every $0<\delta<1$ the same data give the required
radial formula on $\delta\le s<1$.
\lean{PoincareMT.M25.Topology3D.schoenfliesData_of_ball}
\lean{PoincareMT.M25.Topology3D.schoenfliesService_from_main}
\source{Hatcher3M2014}\dcref{Theorem 1.1 and Lemma 1.3, pp. 1--5}
\end{lemma}

\begin{proof}
An ambient diffeomorphism taking the standard
unit sphere to $\psi(S^2\times\{0\})$ supplies a ball chart $B$ with exactly
that boundary.  Choose a sign $\sigma\in\{1,-1\}$ so that
$\psi(q,\sigma s)$ points outward for $s>0$, and let $h:S^2\to S^2$
be the boundary reparametrisation satisfying
$B(h(q))=\psi(q,0)$.  On a radial annulus define
\[
 A((1+s)h(q))=\psi(q,\sigma s).
\]
The radial direction and sphere direction are unique, so this is a genuine
annular chart, and $A=B$ at radius one.  We explain the adjustment that
makes this equality hold on an open overlap.  The transition
$T=B^{-1}\circ A$ is defined near $S^2$, fixes it pointwise, and takes
the exterior side to the exterior side.  Its derivative is the identity
on every tangent plane to $S^2$.  Its normal component is positive:
it is nonnegative by the exterior-side condition and nonzero by
invertibility of $DT$.

Let $\theta:\mathbb R\to[0,1]$ be smooth, zero for $t\le0$ and one
for $t\ge1$.  Near the sphere put
\[
 T_t(x)=x+\theta(t)(T(x)-x).
\]
At a point of the sphere its derivative fixes the tangent plane and
has positive normal component
$1-\theta(t)+\theta(t)\langle x,DT(x)x\rangle$.
Thus $(t,x)\mapsto(t,T_t(x))$ is a local diffeomorphism there.  It is
injective on the compact set $[-1,2]\times S^2$, which it fixes.  There
is consequently one open neighborhood of this compact set on which it
is a chart.  To see why shrinking suffices, a failure of injectivity on
every shrinking neighborhood would give pairs with equal images
converging to two points of the compact set.  Injectivity on that set
makes the limits equal, contradicting local injectivity at the limit.

On the image chart the velocity is
$V(t,y)=\partial_tT_t(T_t^{-1}(y))$.
Multiply it by a smooth compactly supported function equal to one on
the trajectories of a smaller closed neighborhood of the sphere.
The resulting time-dependent vector field has a global smooth flow
$\Phi_t$ on the relevant compact time interval.  Uniqueness of its
integral curves gives $\Phi_t(x)=T_t(x)$ on that smaller neighborhood.
It fixes $S^2$ throughout; no trajectory starting inside or outside
can cross the fixed sphere.  Applying the inverse flow shows that
$\Phi_1$ preserves both the open and closed unit balls exactly.
Consequently $C=B\circ\Phi_1$ has the same filled region as $B$ and
agrees with $A$ on an open neighborhood of $S^2$.

Use $C$ on the inner ball and $A$ on the outer annulus.  Their agreement
on an open overlap makes their union a smooth injective chart $J$ on the ball of
radius two, with smooth inverse.  Its image is the filled interior together
with the whole outward half-collar, and
\[
 J((1+s)h(q))=\psi(q,\sigma s)\quad(\delta\le s<1).
\]
The radial function is explicitly $1+s$, hence positive, strictly increasing,
and less than two on this interval.  This verifies the prescribed collar
data needed when filling a cap; merely identifying its boundary with a
sphere would not give this agreement.
The region $\Omega=B(B_1)$ is bounded because its closure is the image
of the compact closed unit ball.  It and the complementary exterior
are connected, being ambient images of the two standard regions.
The two connected half-collars cannot cross the central sphere, which
gives the stated inward inclusion and exact outer image.
\end{proof}

\begin{lemma}[Orthogonal endpoint for a sphere isotopy]
\label{lem:neck-cap-sphere-isotopy}
For every smooth diffeomorphism $f:S^2\to S^2$ there are an orthogonal
map $A:\mathbb R^3\to\mathbb R^3$ and a jointly smooth family of
diffeomorphisms $I_t$, $t\in\mathbb R$, with $I_0=A|_{S^2}$ and $I_1=f$.
\lean{PoincareMT.M25.Topology3D.diffSphereIsotopyService}
\source{Smale1959}\dcref{Theorems 5--6, pp. 624--626}
\end{lemma}

\begin{proof}
Here is the reduction to compact planar isotopy, retaining its endpoint
convention.  Fix a pole $p$.  The radial extension
$R(x)=|x|f(x/|x|)$ is a diffeomorphism off the origin.  Its derivative
$L=DR(p)$ is invertible and $Lp=f(p)$.  The sphere map
$B(q)=Lq/|Lq|$ therefore has the same first-order behavior as $f$ at
$p$.  The map $g=B^{-1}\circ f$ fixes $p$ and has identity tangent
derivative there.  In stereographic coordinates centered at $p$ its
germ $g_0$ fixes zero and has derivative the identity.

Choose a bump $\chi_r$ supported in $B_{2r}$ and equal to one on
$B_r$.  For sufficiently small $r$ the map
$P(x)=x+\chi_r(x)(g_0(x)-x)$ has $\|DP-I\|\le1/2$ everywhere.
Indeed $\|D\chi_r\|\le C/r$, whereas the identity derivative gives
$\|g_0(x)-x\|\le2r\varepsilon$ and $\|Dg_0-I\|\le\varepsilon$
on $B_{2r}$; hence the derivative error is at most
$(1+2C)\varepsilon$.  The same estimate holds for
$x+\theta(t)(P(x)-x)$.  These maps are bijective: to solve for the
preimage of $y$, apply the contraction theorem to
$x\mapsto y-\theta(t)(P(x)-x)$.  Their derivatives are invertible, so
the parameter-dependent inverse function theorem gives jointly smooth
inverses.  They extend by the identity to sphere diffeomorphisms $C_t$,
with $C_0=\mathrm{id}$ and $C_1=g$ near $p$.  Therefore
$k=C_1^{-1}\circ g$ is the identity near $p$.

Gram--Schmidt applied to the columns of $L$ gives an orthogonal map
$A$ for which $(1-\theta(t))L+\theta(t)A$ stays invertible: in the
orthonormal output basis its matrix is upper triangular with positive
diagonal entries.  Normalize this linear path radially and reverse
its time parameter to obtain $V_0=A|_{S^2}$ and $V_1=B$.
Set $N_t=V_t\circ C_t$.  Then
\[
 N_0=A|_{S^2},\qquad N_1\circ k=f.
\]
Stereographic coordinates turn $k$ into a compactly supported plane
diffeomorphism.  Its compact planar isotopy, extended across the pole,
gives $J_0=k$ and $J_1=\mathrm{id}$.  A common compact support for all
times makes this extension identically the identity on a fixed
neighborhood of the omitted pole, so joint smoothness holds there too.
Thus
\[
 I_t=N_t\circ J_{1-t}\quad(0\le t\le1)
\]
is a smooth family of diffeomorphisms from $A|_{S^2}$ to $f$.  This permits
either orientation and explains why an isotopy from the identity alone
would not suffice for arbitrary boundary parametrisations.
\end{proof}

\subsection{The compact planar isotopy}
\label{sec:neck-cap-planar-isotopy}

We now supply the planar input used above.  If $h$ is a smooth
diffeomorphism of $\mathbb R^2$ equal to the identity outside a compact
set, there is a jointly smooth family $M_t$, with jointly smooth inverses,
such that $M_t=h$ for $t\le0$, $M_t=\mathrm{id}$ for $t\ge1$, and
all $M_t$ are the identity outside one compact set.  The argument is the
relative smooth version of Munkres' Theorem 1.3 and Lemmas 1.1--1.2,
pp.~193--195 \cite[Theorem 1.3 and Lemmas 1.1--1.2, pp.~193--195]{Munkres1960}.
His statement is a regular isotopy at finite
differentiability; the smooth rounding and relative constructions below
provide the joint smoothness required here.

First move the support above the axis.  If it lies in the ball of radius
$R$, choose $v=(0,-R-1)$ and a smooth function $\alpha$ equal to zero
for $t\le1/3$ and one for $t\ge2/3$, with values in $[0,1]$.  Set
\[
 N_t(x)=h(x+\alpha(t)v)-\alpha(t)v,\qquad g=N_1.
\]
Both this family and its inverse are smooth and are supported in the
compact union of the translates of the original support over $[0,1]$.
The map $g$ is the identity for $y<1$ and for $y\ge2R+2$.
Choose $b=4R+4$ and slide in the opposite direction:
\[
 F_t(x,y)=g(x,y+b\alpha(t))-(0,b\alpha(t)).
\]
It starts at $g$ and ends at a map which is the identity on the upper
half-plane.  It need not end at the identity below the axis.  Near both
time ends it fixes one uniform strip about the axis.  Its moving axis
$C_t(u)=F_t(u,0)$ is an embedded regular line, equal to $(u,0)$ outside
one compact interval, and equal to that line near the two time ends.

To correct this moving axis, take a uniform thin normal tube around the
family $C_t$ on a compact time interval.  Such a tube exists because
the tangent never vanishes, local inverse charts cover the compact moving
part, and pairs outside those charts have a positive separation; the
fixed straight tails use the standard tube.  Inscribe a sufficiently fine
affine-grid polygon in $C_t$, retaining those tails.  Compact uniform
continuity of position and tangent makes the polygon simple and its
chords close to the corresponding tangents.  Round each vertex by a
smooth convex replacement of absolute value, equal to absolute value
off a small interval and with derivative bounded by one.  The rounded
line $\gamma_t$ is jointly smooth and its normal-tube coordinates
$(\xi_t(u),\eta_t(u))$ satisfy
\[
 \partial_u\xi_t(u)>0,\qquad |\eta_t(u)|<w/2,
\]
where $w$ is the tube width.  The same rounded samples will be used in
both transports below.  On the fixed tails and on stationary time slices
they are exactly the original straight line.

Here is the finite polygon argument behind the second transport.
An admissible interior vertex is one whose adjacent triangle meets the
arc only in its two adjacent edges.  Moving it to the midpoint of its
neighbors inside that triangle preserves simplicity and makes it
removable.  An admissible vertex persists on a small parameter window.
Finitely many such windows cover the time interval; choices at successive
windows may be made equal or nonadjacent.  Supported pushes therefore
preserve the straight vertex needed in the neighboring window.  Near
each junction first use induction on the smaller number of vertices to
make the whole arc straight.  On each remaining interval delete its
straight vertex, apply the same induction relative to the straight
endpoint windows, and reinsert the vertex.  The interval homotopies
agree on those fixed windows and paste smoothly.  This gives a smooth
two-parameter polygon family with fixed endpoints and time tails,
ending in the straight segment.  The admissible-vertex choice and this
relative induction are Munkres, Lemmas 2.3--2.7, pp.~195--197.

Rounding that finite polygon homotopy with one sufficiently small
profile gives an isotopy of embedded smooth lines.  Subdivide its homotopy
parameter into finitely many intervals on which each line is a small
normal graph over the line at the interval's initial parameter.
Supported graph transports, composed successively, give an ambient
map $J_t$ taking $\gamma_t$ to the axis, with $J_t=\mathrm{id}$ at
the time ends.  There is also a relative ambient map $I_t$ taking
the original $C_t$ to $\gamma_t$.  Indeed $\xi_t$ is increasing and
is the identity on the tails, so it is a diffeomorphism of the line
with smooth parameter-dependent inverse.  In tube coordinates
$\gamma_t$ is thus the graph
$\eta_t\circ\xi_t^{-1}$.  Write its height as $a_t(u)$, with
$|a_t|<A<w$.  Choose a normal cutoff $\beta$ supported inside the
remaining tube width, with $\beta(0)=1$ and $|\beta'|\le B$,
and an integer $n>BA$.  The $k$th small slide is
\[
 (u,y)\longmapsto
 \left(u,y+\beta\bigl(y-ka_t(u)/n\bigr)a_t(u)/n\right),
 \qquad 0\le k<n.
\]
Its normal derivative is positive because $B|a_t|/n<1$, and its
normal coordinate is the identity outside a bounded interval, hence
is bijective.  It takes height $ka_t/n$ to $(k+1)a_t/n$.
Conjugate these maps by the tube chart and extend by the identity;
their supports lie compactly inside the chart.  Their composition
is the required graph transport.  It is the identity whenever the graph
is zero, and has one compact support because the parameter range and
nonstraight part are compact.  Hence
$G_t=J_t\circ I_t$ takes $C_t$ into the axis and is the identity at
both time ends.  This develops the normal-slide step of Munkres,
Section~2.8, p.~197, using smooth rounded samples throughout.

It remains to fix a whole strip, not merely the axis setwise.
The restriction of $G_t\circ F_t$ to the axis is increasing: its
derivative is nonzero and it is the identity on the tails.  Correct this
one-dimensional map by a supported ambient interpolation so that the
axis is fixed pointwise.  The resulting derivative along the axis is
\[
 \begin{pmatrix}1&a(t,u)\\0&c(t,u)\end{pmatrix},\qquad c(t,u)>0.
\]
The positivity follows from orientation preservation, which holds since
the plane map is the identity outside a compact set.  Realize this
derivative by a supported positive normal scaling and horizontal shear,
and compose with its inverse.  Now the map fixes the axis with identity
derivative.  On a uniformly thin strip its displacement is $o(|y|)$
and its derivative error is $o(1)$, uniformly in $t$ and $u$ on the
compact moving part.  Multiplying that displacement by a strip cutoff
therefore gives a global map whose derivative differs from the identity
by less than $1/2$.  The contraction theorem makes it a diffeomorphism;
composing with its inverse corrects the original map on a narrower strip.
All corrections vanish at the time ends.  We have constructed $A_t$
with one compact support and
$A_tG_tF_t=\mathrm{id}$ on $|y|<\delta$, for some $\delta>0$.

Put $H_t=A_tG_tF_t$ on $y\ge0$ and $H_t=\mathrm{id}$ on $y\le0$.
These formulas agree on an open strip, so they and their inverses paste
smoothly.  The corrected plane map preserves both half-planes, as it fixes
the strip and is the identity near infinity.  Thus $H_t$ is a global
diffeomorphism, with $H_0=g$ and $H_1=\mathrm{id}$.  Finally the smooth
composition
\[
 M_t=N_{2t}\circ g^{-1}\circ H_{2t-1}
\]
is $h$ for $t\le0$ and the identity for $t\ge1$.  Each factor and
its inverse fixes the complement of a common compact set.  This proves
the support, endpoint and smoothness assertions needed to extend the
planar family across the stereographic pole.

\section{Local and global topological alternatives}
\label{sec:m25}

The local theorem fixes one threshold before the manifold, metric, or cover.
The current proof takes the minimum of eight positive thresholds: the
separating-chain tube, its whole-cover specialization, uniform separation
labels, the nonseparating circle-bundle construction, the separating local
branch, the nonseparating local branch, the cap-core frontier construction,
and the finite capped-chain construction.  Consequently
$0<\varepsilon_0\le 1/200$, and every cover with $\varepsilon\le\varepsilon_0$
inherits all four conclusions at the same $\varepsilon$.

For a neck-only cover, the separating case supplies a balanced tube whose
source necks lie in the original cover, whose selected centres lie in $X$,
whose carrier contains $X$, and whose necks are all separating.  Under the
additional whole-cover hypothesis, uniform separation labels give either this
tube with carrier exactly $M$, or the sphere bundle over a circle described
above.  For a connected neck--cap cover, the local conclusion is one of six
compatible regions: two caps forming a closed component, a double-capped tube,
a single cap, a capped tube, a tube, or a sphere bundle.  Compatibility means
that each displayed neck/cap has the cover's $\varepsilon$ and every cap
constant is bounded by the cover constant.

The whole-cover corollary turns the six local cases into four global cases:
\begin{enumerate}
\item a whole closed component modelled on $S^3$, $\mathbb{RP}^3$, or
      $\mathbb{RP}^3\#\mathbb{RP}^3$;
\item a noncompact cap or capped tube whose model is Euclidean or punctured
      projective;
\item a whole balanced tube; or
\item a whole smooth sphere bundle over the circle.
\end{enumerate}
The implication uses only $X=M$: a local containment $X\subset Y$ becomes
$M\subset Y$, and hence $Y=M$; no new smallness estimate is introduced.

\section{Global alternatives for ancient solutions}
\label{sec:m26}

Let $K$ be a connected, second-countable ancient three-dimensional
$\kappa$-solution, with time domain $t\le0$.  A strong evolving neck at time
$t$ has a terminal $\varepsilon$-neck, the same flow connection, and a
backward slab of duration $d$ satisfying
\[
 dR(p,t)=1,
\]
on which the rescaled metric remains close to the round evolving cylinder.
A strong cap is a cap certificate on the time-$t$ slice with the actual flow
connection and uniform scalar, diameter, volume, gradient, and evolution
estimates.  Strong tubes and strong capped tubes require strong-neck centres
throughout their tube carrier; a double-capped tube retains two disjoint
closed cores and both cap attachments.

The preceding analytic inputs include curvature and scalar evolution, the
pointwise curvature bound by later scalar curvature, scalar normalization,
fixed-carrier blow-up setup, two- and three-dimensional asymptotic
classification, universal noncollapsing, normalized compactness, model
refinement, and the whole-cover conclusion above.  These hypotheses are
quantified before $K$ is chosen.  The proof preserves that order.
In particular the earlier curvature estimate means
$|\operatorname{Rm}|(x,t)\le R(x,b)$ for every $t\le b\le0$ at the
same point, and normalization means the actual flow
$R(p,b)g(b+s/R(p,b))$, with $R(p,b)>0$.  Universal noncollapsing
provides one $\kappa_0>0$ for every nonround solution before normalized
compactness is applied.  The compactness conclusion retains a single
carrier and convergence maps through time zero, as in
Theorems~\ref{thm:ancient-universal} and~\ref{thm:ancient-compactness}.
The classifications and fixed-time model identifications are those of
Theorems~\ref{thm:surface-ancient}, \ref{thm:three-shrinkers}, and
\ref{thm:ancient-model-identifications}; the fixed-carrier rescaling setup
is Proposition~\ref{prop:ancient-sequence}.  These are explicit inputs
to the arguments below.

For the noncompact branch, the original-flow curvature trichotomy leaves a
sphere-line model, a positive-curvature model, or a twisted quotient model;
the projective-plane-line model is excluded by the no embedded trivial-normal
projective-plane hypothesis.  The sphere-line model has strong neck coverage
everywhere and hence gives a whole strong tube.  In the positive model, a
compact soul construction provides a cap; maximal outgoing chains cover its
exterior by strong necks, yielding a strong capped tube.  The twisted model
has the same capped-tube conclusion with the punctured-projective cap.  Taking
the minimum of the three accuracy thresholds and the maximum of their cap
constants gives the quantifier order
\[
 \exists\varepsilon_2>0\ \forall\varepsilon\in(0,\varepsilon_2]\quad
 \exists C_0>0\ \forall K.
\]
Here $K$ ranges over the noncompact solutions with the stated projective-plane
exclusion; $C_0$ may depend on $\varepsilon$, but neither constant depends on
the particular solution or underlying manifold.

For a compact slice, normalized compactness gives a threshold $\varepsilon_3$
and constant $C_1$.  The alternatives are: a round ancient quotient, a
compact small-diameter positive component, or a double-capped strong tube.
The round branch carries a finite spherical-space-form quotient, not only the
simply connected sphere.  This is a zero-time slice
statement; it does not yet assert a canonical neighborhood at every ancient
time.

\section{Pointwise canonical neighborhoods}
\label{sec:m27}

For this refinement we use the preceding compact and noncompact
alternatives, the separating-tube and whole-cover results, scalar and
curvature evolution, scalar normalization, universal noncollapsing,
normalized compactness, and the two-dimensional and model-refinement
inputs.  The construction below uses these statements directly.  In
particular it does not assume the additional blow-up setup, classified
asymptotic limit, or Ricci-evolution service in the preceding global
theorem's full input bundle.  Its compact two-core conclusion is
constructed from the refinement's own inputs.

We now add the uniform scalar derivative bounds.  There is a constant
$0\le B<C$ such
that, for every $t\le0$ and $x\in M$,
\[
 |\nabla R|(x,t)\le B R(x,t)^{3/2},\qquad
 |\partial_tR(x,t)|\le B R(x,t)^2,
\]
with the time derivative interpreted as a derivative within $(-\infty,0]$.
The proof normalizes at $(x,t)$, applies the uniform curvature-jet estimate,
contracts the first two curvature derivatives to scalar estimates, and uses
the Ricci-flow scalar equation.  Round solutions are handled separately and
then combined with the nonround estimate by increasing the common constant.
For clarity, normalized compactness supplies one scalar bound on a unit
ball for each fixed noncollapsing constant.  The bound by later scalar
curvature controls the curvature there backwards in time, so the local
derivative estimates give constants $D_1,D_2$ for the first two curvature
derivatives at the normalized centre.  Their scalar contractions have
bounds $9D_1$ and $27D_2$.  Nonnegative Ricci curvature and $R=1$ give
$|\operatorname{Ric}|^2\le1$ there, hence
$9D_1+27D_2+2$ bounds both scalar expressions.  Use the universal
$\kappa_0$ for nonround solutions and enlarge this constant for the
explicit round calculation.  This explains the uniform normalized jet
input without assuming the scalar derivative conclusion itself.

The powers of scalar curvature follow directly from parabolic scaling.
At a fixed $(x,t)$ let $\lambda=R(x,t)>0$ and set
$\widetilde g(s)=\lambda g(t+s/\lambda)$ for $s\le0$.
Then $\widetilde R(x,0)=1$,
\[
 |\widetilde\nabla\widetilde R|(x,0)
       =\lambda^{-3/2}|\nabla R|(x,t),\qquad
 \partial_s\widetilde R(x,0)=\lambda^{-2}\partial_tR(x,t).
\]
The first identity uses $\widetilde R=\lambda^{-1}R$ and the factor
$\lambda^{-1/2}$ in the norm of a covector.  The second uses the same scalar
factor and $dt/ds=\lambda^{-1}$.  Applying one uniform normalized jet bound
and multiplying by the indicated powers proves both estimates with a
constant independent of $(x,t)$ and of the solution.  At $t=0$ the chain
rule is within the closed ancient time interval, so no future extension of
the solution is required.  Increasing the final neighborhood constant
strictly above this uniform bound gives the displayed slack $B<C$.

More precisely, there exists $\bar\varepsilon>0$ such that, for every
$0<\varepsilon<\bar\varepsilon$, there is $C>0$, independent of the solution
and its underlying manifold, giving the derivative bounds and the following
seven alternatives at time zero:
\begin{enumerate}
\item a round spherical quotient;
\item a compact positive-curvature geometry with topology $S^3$ or
      $\mathbb{RP}^3$, scale-invariant diameter, volume, and sectional bounds;
\item a double-capped strong tube with two retained cap cores;
\item a capped Euclidean tube with cap-core-or-strong-neck coverage;
\item a capped twisted quotient with the same coverage;
\item a sphere-line product with a strong tube; or
\item the projective-plane-line product.
\end{enumerate}
In the compact positive branch, put $D=\operatorname{diam}_{g(0)}M$ and
$V=\operatorname{Vol}_{g(0)}M$.  The estimates are, for every scalar basepoint
$x\in M$,
\[
 C^{-1/2}R(x,0)^{-1/2}<D<CR(x,0)^{-1/2},\qquad
 C^{-1}R(x,0)^{-3/2}<V<CR(x,0)^{-3/2}.
\]
For every independently chosen point $y$ and every orthonormal pair $a,b$ at
$y$, one has
\[
 C^{-1}R(x,0)<\operatorname{sec}_{g(0)}(a,b)<CR(x,0).
\]
The independence of $x$ and $y$ is part of the conclusion.  The double-capped
branch also has positive sectional curvature and topology $S^3$ or
$\mathbb{RP}^3$; every point belongs to one of its two open cap cores or is
the centre of a strong evolving neck.  The Euclidean capped branch has
positive sectional curvature and a smooth identification with $\mathbb R^3$.
The quotient capped branch instead retains its twisted sphere-line flow
model.  These last two branches have the analogous one-core-or-neck coverage.

The compact-positive proof first handles bounded scalar-scale diameter by
normalized compactness and positivity.  If diameter is unbounded, select
points with neither a strong neck nor a controlled strong-collar cap, normalize,
and pass to a normalized ancient limit.  The limit is noncompact.  Its cap or product model
gives a strictly finer strong-neck cover outside the open cap core, while
full-domain neck stability transports the cap and its compact collar back to
the chosen sequence, a contradiction.  The remaining whole cover has strong
collars: every point of a cap carrier outside its open core is a strong-neck
centre.  This is stronger than saying that the complement of the cap carriers
has strong necks.

There is a further quantitative split before the two caps are selected.
Let $c$ be the common cap bound supplied by the collar cover and $d$ the
scalar-scale diameter bound in the first alternative.  The encounter
construction has the two-cap diameter bound $d_2(c)=2c/c^{-1/2}+1$.
Apply the bounded-diameter
geometry theorem with $b=\max\{d,d_2(c)\}$, obtaining a geometry constant
$C_g$.  The final constant can be $\max\{C_g,c\}$.  If the diameter is less
than $d_2(c)R(x,0)^{-1/2}$ for every $x$, this is again the compact-positive
branch.  Otherwise there is an actual point $p$ such that
\[
 d_2(c)R(p,0)^{-1/2}\le\operatorname{diam}_{g(0)}M.
\]
This witness, together with compactness and the strong collars, is the
large-diameter input to the two-cap selection theorem.  Positive sectional
curvature supplies the exclusion of a two-sided embedded projective plane;
no extra orientability premise is inserted.

The selection produces two actual caps $A,B$ from the cover and a finite
balanced connecting chain.  Its obstruction statement says that a point
outside both open cores which is not a strong-neck centre would give a cap
$D$ whose entire closed core is contained in that finite chain.  The
closed-core noncontainment lemma for a quarter-capturing chain rules this
out.  Thus every point outside $A$ and $B$'s cores is a strong-neck centre.
The matching construction now makes a strong double-capped tube while
preserving the equalities of its two open cores with the originally selected
cores.  These equalities transport the just-proved coverage statement to the
final tube.  Replacing the common constant by $\max\{C_g,c\}$ weakens only
the quantitative bounds; it does not change either core or the coverage.
This is the retained two-core double-capped branch.

We spell out the topological obstruction in this last step.  Suppose a cap's
closed core were contained in a finite quarter-capturing chain of
equal-parameter necks.  Its boundary sphere is then contained in the chain
union.  The sphere-capture estimate places that entire sphere in one selected
neck.  It uses the inner three-quarter slab, scale comparison at a common
closure point, and the fact that the diameter of a central sphere is at most
$2\pi$ times its scale: for $\varepsilon\le1/1000$, the resulting distance
upper bound is smaller than the $.99r\varepsilon^{-1}$ cost of leaving the
selected neck.  If there is no inner-slab contact, quarter capture and the
positive-end exclusion instead force the same containment at an adjacent
neck.  Thus this is containment of the entire sphere, not only its centre.

A graph transport supported compactly inside that neck moves its central
sphere to the cap boundary.  Pulling the core back by this ambient
homeomorphism preserves compactness, nonempty interior, and containment in
the chain union, and makes its frontier the selected central sphere.
Successive graph transports along the finite chain then move that sphere
to the middle sphere of the chain's open-cylinder model.  All transports
preserve the exact chain union.  We would therefore have a compact subset
of $S^2\times(0,1)$ with nonempty interior and frontier exactly
$S^2\times\{1/2\}$.  On either connected half-cylinder its intersection is
both open and closed, because there are no frontier points there.  Nonempty
interior meets one of the halves, so that entire half must be included.
But a compact subset of the open cylinder has its interval coordinate
bounded away from both zero and one.  Containing a whole half-cylinder is
impossible.  This contradiction proves the closed-core noncontainment used
to establish strong-neck coverage.

The all-time corollary chooses a separate $\varepsilon'>0$, then for each
$0<\varepsilon\le\varepsilon'$ a constant $C>0$ before the solution.  Outside the
projective-plane-line certificate, every $t\le0$ and every $x\in M$ lies in a
strong neck, a quantitative cap, a compact component, or a round component.
The projective-plane-line exception is explicit and is not silently absorbed
into the cap branch.  The all-time conclusion supplies the canonical
alternatives used by later surgery constructions; the preceding global
classification concerns the distinguished time-zero slice.

\section{Cross-chapter contracts and limits}
\label{sec:neck-cap-contracts}

The ancient-models chapter must provide the fixed-carrier limits, asymptotic
classification, noncollapsing and compactness used by the ancient alternatives,
with the precise meanings specified in Section~\ref{sec:m26}.
Together with the curvature and scalar evolution equations from the
Ricci-flow chapter, these give the normalized scalar jets used in
Section~\ref{sec:m27}.  The singular-limits chapter consumes the local neck/cap
certificates but must separately prove transfer from an ancient limit to a
high-curvature point in the original surgery flow; the ancient alternatives do not perform
that blow-up transfer.  Its terminal cap normalization also uses the
full-carrier estimate \eqref{eq:neck-cap-full-scalar}, with its linear
error and unchanged neck scale.  The surgery chapter consumes the actual smooth ball
and sphere gluing, cap cores, collar attachments, and the pointwise canonical
alternative; it must construct the standard cap and metric interpolation.
The sphere-reduction chapter consumes the closed-component certificates and
the collared sphere gluing, then proves the finite connected-sum reduction.

The ancient classification, normalized compactness, and evolution inputs
remain hypotheses until supplied by their earlier chapters.  Their precise
integration status, the source revisions, and the outstanding review points
are recorded in the companion review artifact.

\section*{Sources}

The neck and cap definitions and the corrected Appendix A construction are
Morgan--Tian, Definitions 9.72, 9.77, and 9.81, and Definitions and Lemmas
A.12--A.18, Proposition A.19, Lemma A.20, Proposition A.21 with
Claims A.22--A.24, and Proposition A.25, printed pp.~230--235 and
504--514.  The ancient alternatives are Morgan--Tian, Corollary 9.88,
Theorem 9.89 and Claim 9.90, Theorem 9.93 and Corollary 9.94, printed
pp.~239--243.  Sphere extension uses Hatcher, \emph{Notes on Basic 3-Manifold
Topology}, Theorem 1.1, pp.~1--5; the sphere isotopy input uses Smale,
Theorems 5--6, pp.~624--626.  The compact planar construction follows
Munkres, Theorem 1.3 and Lemmas 1.1--1.2, pp.~193--195, and Lemmas
2.1--2.7 with Section~2.8, pp.~195--197; its jointly smooth relative
version is developed in Section~\ref{sec:neck-cap-planar-isotopy}.

\lean{PoincareMT.m25NeckCapTopology}
\lean{PoincareMT.m26CanonicalNeighborhoods}
\lean{PoincareMT.m27KappaAlternatives}
\source{MT2007}
\dcref{Appendix A, pp. 504--514; Corollary 9.88--Corollary 9.94, pp. 239--243}
\cite[Appendix A, pp.~504--514; pp.~239--243]{MT2007}
\cite[Theorem 1.1, pp.~1--5]{Hatcher3M2014}
\cite[Theorems 5--6, pp.~624--626]{Smale1959}
